% Ad Familiares 10.34 --- headnote
% ad-familiares-10-34, 18 May 43 BC, camp at Pons Argenteus

\textit{M. Aemilius Lepidus to Cicero, written from camp
at the river Argenteus around 18 May 43 BC --- Perseus
dateline \textit{Scr. in castris ad Pontem Argenteum
circ. xv K. Iun. a. 711 (43)}. Lepidus writes as
\textit{imperator iterum} and \textit{pontifex maximus}
(the Latin salutation transmitted in Perseus reads PONL
for what is plainly PONT., and is preserved as
transmitted), reporting in the bureaucratic dispatch
style of a provincial commander to a senior senator. The
position he describes is the camp at Pons Argenteus on
the river Argens in Gallia Narbonensis, where the
remnant of Antony's army --- driven west out of Italy
after the defeat at Mutina in April --- has finally
made contact with Lepidus's seven legions. Antony's
forces, Lepidus reports, are visibly dwindling: men
desert daily, Silanus and Culleo have come over
already, and the only thing keeping Antony's army in
the field is Ventidius's four legions and a cavalry
arm of more than five thousand.}

\textit{The closing sentence --- ``As far as this war
is concerned, we shall not fail the Senate or the
commonwealth'' --- is the assurance Cicero and the
Senate had been demanding of Lepidus for weeks, and it
is the assurance Lepidus would publicly repudiate
twelve days later in the dispatch that survives as
10.35. By the time he writes this letter the choice
between Antony and the Senate has narrowed to days;
the troops on both sides --- veterans of Caesar, many
of them with personal ties across the line of battle
--- have begun fraternizing, and the chain of command
on Lepidus's side is already slipping. The cool
official tone, the careful itemization of Antony's
weakness, the catalogue of defectors named one by one
--- all of it reads, in the light of what came on 29
May, as a man writing for the record he expects to be
held against him.}
